\chapter*{Úvod}
\addcontentsline{toc}{chapter}{Úvod}

Bitcoin umožňuje uživatelům spravovat finanční prostředky bez prostředníků,
ale uživatel za to platí tím, že se musí starat o své privátní klíče. Ztráta nebo
zcizení klíče znamená nevratnou ztrátu prostředků. Kvůli tomu se používají
hardwarové peněženky, které uchovávají klíče na izolovaném zařízení,
a~multisignature schémata, kde je k~provedení transakce potřeba souhlas více
klíčů. Pokud uživatel obě metody zkombinuje, výrazně tím sníží riziko ztráty
i~krádeže.

Na desktopu tuto kombinaci pokročilých funkcí nabízí například Sparrow
Wallet~\cite{SparrowWallet} nebo Electrum~\cite{Electrum}. Na Androidu je ale
situace horší, většina peněženek buď multisig nepodporuje, nebo jej
implementuje jen omezeně, bez přímé integrace s~hardwarovou peněženkou.
Podobně je na tom coin control, tedy ruční výběr konkrétních UTXO při tvorbě
transakcí.

Cílem práce je navrhnout a~implementovat mobilní aplikaci pro Android, která
tyto nedostatky řeší. Aplikace podporuje multisignature transakce ve schématu
M-of-N, integraci s~hardwarovou peněženkou Trezor a~manuální výběr UTXO.
Server přitom nikdy nepracuje s~privátními klíči, transakce se podepisují
výhradně na Trezoru přes rozhraní Trezor Connect~\cite{TrezorConnect}.

Pro distribuci částečně podepsaných transakcí mezi cosignery používáme formát
PSBT (BIP-174)~\cite{BIP174}. Backend tvoří soustava mikroslužeb v~Kotlinu
s~frameworkem Ktor~\cite{KtorFramework}, frontend je nativní Android aplikace
v~Jetpack Compose~\cite{JetpackCompose}.

Práce je členěna následovně. Kapitola~1 pokrývá teoretické základy Bitcoinu
a~srovnání existujících peněženek. Kapitola~2 specifikuje požadavky formou
user stories a~use cases. V~kapitole~3 popisujeme návrh architektury,
technologická rozhodnutí a~vybrané implementační detaily včetně technických
výzev, na které jsme narazili. Kapitola~4 popisuje testování na testnetové
síti a~kapitola~5 obsahuje uživatelskou dokumentaci.
